% gamma-v1 — Paper A adapted for Digital Humanities Quarterly (DHQ).
% Base content: paper/v7 (firewall-sourced). A DH-framed, epistemology-first sibling of the
% alpha (TACL) draft — NOT the same manuscript (see the alpha/gamma coordination note in
% docs/planning/publication-h2-2026/pathways/gamma/v1.md). Article class for portability;
% DHQ uses an XML authoring workflow at acceptance. Single-blind: author retained.
\documentclass[11pt]{article}
\usepackage[utf8]{inputenc}
\usepackage[T1]{fontenc}
\usepackage{lmodern}
\usepackage[margin=1in]{geometry}
\usepackage{microtype}
\usepackage{amsmath,amssymb}
\usepackage{enumitem}
\usepackage{booktabs}
\usepackage{xcolor}
\usepackage[hidelinks]{hyperref}
\usepackage{natbib}
\usepackage{titlesec}
\titleformat{\section}{\normalfont\large\bfseries}{\thesection}{0.6em}{}

\title{\textbf{What the Statistics Can and Cannot Say About the Voynich Manuscript:\\
A Firewall-Disciplined, Evidence-Graded Constraint Study}}
\author{Tim Walsh \\ \small DireLabs, LLC}
\date{}

\begin{document}
\maketitle

\begin{abstract}
\noindent The Voynich Manuscript (Beinecke MS~408) has drawn a century of confident, mutually
incompatible ``solutions,'' a pattern enabled by a specific epistemic predicament: with no
ground truth against which to check a reading, ``it looks like language / a cipher / a number
system'' is nearly unfalsifiable. We take the opposite stance. Rather than claim what the book
says, we ask, for each interpretable dimension, \emph{how far along it sits}, and we require every
such claim to carry an explicit A--D evidence grade produced by deterministic, reproducible code.
We confirm the manuscript's anomalously low character entropy and its two-system (Currier~A/B)
structure; we show \emph{why} purely distributional statistics cannot localize meaning
(structured-but-meaningless processes reproduce them); and we turn a mid-level syntactic signature
into a cryptanalytic discriminator that robustly excludes word-order-preserving ciphers of real
prose while leaving verbose-homophonic ciphers inconclusive. The contribution is a disciplined way
to \emph{shrink the space of viable hypotheses} about an undeciphered artifact --- and an honest
account of where distributional methods stop. We make no decipherment claim.
\end{abstract}

\section{Introduction: a discipline of not-knowing}
For more than a century, readers have announced that the Voynich Manuscript is solved --- as
medieval Latin shorthand, as an artificial language, as Hebrew, as proto-Romance, most recently by
artificial intelligence --- and each time the wider community has shown the reading does not hold
\citep{bowern2021linguistics}. This is not merely a run of bad luck. It is the predictable
consequence of studying an object with \emph{no recoverable ground truth}: when nothing can be
checked against a key, any sufficiently clever reading ``comes out looking like'' sense, and there
is no principled way to tell an insight from a coincidence.

This paper adopts the discipline that predicament demands. We do not offer a reading. We treat the
manuscript as an object to \emph{constrain} --- to ask which broad classes of explanation remain
consistent with the evidence and which can be set aside --- and we hold ourselves to two rules
that the field's history of over-claiming motivates: every quantitative claim is produced by
deterministic, versioned code (so it can be re-run and audited, never merely asserted), and every
claim wears an explicit grade from A (robust) to D (speculative). The result is not an answer but
a \emph{smaller, better-defended space of possible answers}, and --- we will argue --- a clearer
view of exactly where quantitative methods lose their grip.

\section{The manuscript and the evidence}
MS~408 is an early-fifteenth-century illustrated codex whose script no one can read. We work from
the standard Zandbergen--Landini transliteration of the text \citep{zandbergen_voynichnu} and
stratify throughout by the two writing systems Currier identified \citep{currier1976} and by
scribal hand \citep{fagindavis2020hands}; our measures do not depend on any assumed sound values
for the glyphs. Two facts bound everything that follows: roughly an eighth of the leaves are lost,
and the vocabulary does not saturate --- new words keep appearing --- so every count is a lower
bound on what the complete book contained.

\section{How we know: method as epistemology}
Our method is mostly a set of rules for not fooling ourselves.
\begin{itemize}[leftmargin=1.2em,itemsep=1pt]
  \item \textbf{Validate before claiming.} No measurement is trusted on the manuscript until it
  can tell real language apart from deliberately \emph{meaningless} imitations --- texts generated
  to mimic surface structure while carrying no content --- and from ciphers of known texts. A
  measure that cannot make that distinction on material where we know the answer is not allowed to
  pronounce on the manuscript, where we do not.
  \item \textbf{Numbers only from code.} Every figure we report is written by versioned software
  that records how it was produced. Nothing is estimated from memory or asserted by hand.
  \item \textbf{Attack every strong claim.} Before a strong finding stands, an independent reader
  --- given the evidence but not our story about it --- is asked to make the best possible case
  that we are wrong, and that case is preserved. More than a dozen of our own first conclusions
  did not survive this step; we report the survivors and the casualties alike.
  \item \textbf{Grade everything.} Each claim below carries an A--D grade so a reader can see, at a
  glance, how much weight it can bear.
\end{itemize}

\section{What the statistics say}
\textbf{An anomalous, two-system script (grade A/B).} The manuscript's letter-to-letter
predictability is far higher than any natural-language comparison we tested --- its conditional
character entropy sits well below that of medieval Latin, Italian, German, or Hebrew. Currier's two
``languages'' are genuinely distinct as statistical systems, and the writing is built from a dense
web of near-identical word-forms, more like a paradigm of variations on a theme than an ordinary
vocabulary. (Exact values and comparisons are given in the appendix; here we keep to what they
mean.)

\textbf{No measurement localizes meaning (grade C).} The central negative, and the paper's
epistemic core: across every statistic we can compute, \emph{a deliberately meaningless imitation
scores as ``meaningful'' at least as strongly as real text does}, and the manuscript itself sits
at or below both. There is, on our evidence, no quantitative dial that reads out ``this carries
meaning.'' This is why the honest question is not \emph{what does it say} but \emph{how far along
each measurable dimension does it sit} --- and it marks the exact boundary past which numbers stop
being able to help.

\textbf{Not a cipher of ordinary prose --- with one live exception (grade B/C).} If the manuscript
were a cipher that preserved the word order of some real language, it would inherit that language's
strong sentence-level grammar. It does not: its word-level syntactic structure is weak, whereas a
$1{:}1$ substitution, a vowel-dropping abjad, and a nomenclator applied to Latin all show the
strong structure of their source --- and this holds across typologically diverse languages,
including a native abjad (Hebrew). Once the comparison is made fairly, the gap is large and stable.
\emph{So the manuscript is not a word-order-preserving cipher of real prose.} We had first claimed
more --- that \emph{no} cipher of real prose could fit --- but engaging a
concurrently-published, hand-built cipher designed to imitate the manuscript
\citep{greshko2025naibbe} forced us to withdraw the stronger claim: the word-order evidence we had
leaned on turns out to be confounded by re-spacing and by the use of many interchangeable symbols,
not a clean signal. A verbose, ``many-symbols-per-letter'' cipher therefore remains a live
possibility, as does a rule-driven generative process; an independent recent analysis reaches a
neighbouring, cipher-leaning conclusion \citep{parisel2026layered}.

\section{Where distributional methods stop}
The payoff of the discipline is not a headline but a boundary. We can say, with graded confidence,
what the manuscript's statistics are \emph{like}; we can rule out whole families of explanation; and
we can show that the one question everyone wants answered --- does it mean anything, and if so what
--- is, on present evidence, \emph{provably underdetermined by the statistics}, because meaningless
processes reproduce the same numbers. That is a genuinely useful thing to know: it tells future work
that progress on meaning will have to come from outside the distributional toolkit --- from the
images, the material object, the historical record --- not from another pass over the character
counts.

\section{In conversation with manuscript studies}
Our stance is continuous with the recent turn in computational work on the manuscript that tests
\emph{structure} without claiming to read the text --- for example, latent-semantic analysis of
topical coherence \citep{layfield2026lsa} and the broader linguistic synthesis of
\citet{bowern2021linguistics}, both of which treat the manuscript as something to characterize
rather than crack. Where earlier computational readings sought a solution, this line asks what can
be responsibly said without one; we contribute a reproducible, self-auditing way to do so, and a
tool others can run.

\section{Limitations}
The evidence is a single manuscript, and an incomplete one, so several findings are
power-limited; two of our word-level syntactic measures are ``soft'' (their confidence intervals
cross zero) and we do not lean on them beyond the large cipher-versus-manuscript gap; the
word-order statistic is confounded as noted; and our analysis is based on one transliteration
system with a second used only as a spot-check. The grades above reflect these limits honestly.

\section{Conclusion}
Treating the Voynich Manuscript as an object to constrain rather than to solve turns a century of
irreproducible ``solutions'' into a set of reproducible, gradeable statements: an anomalous
two-system script, a meaning that no statistic can localize, and a cipher hypothesis narrowed but
not closed. The transferable contribution is the discipline itself --- validate before claiming,
numbers only from code, attack every strong claim, grade everything --- offered to a field that has
long needed it.

\section*{Data and reproducibility}
The transliteration is the public Zandbergen--Landini edition; all statistics are produced by
deterministic, versioned code, and the discriminator tool, its reference data, and the archived
adversarial-review briefs are released open-source. Numerical appendix and code links accompany the
article.

\bibliographystyle{plainnat}
\bibliography{refs}
\end{document}
